\documentclass[11pt]{article}
\usepackage[margin=1in]{geometry}
\usepackage{amsmath,amssymb,booktabs,hyperref,graphicx}
\title{Independent Replication --- QAOA Reachability Deficits\\
\large arXiv:1906.11259 (Akshay, Philathong, Morales, Biamonte, 2019)}
\author{Ollie (Rick Stevens, Argonne National Laboratory)}
\date{2026-07-03 (backfill 2026-07-06)}
\begin{document}
\maketitle

\section{Verdict}
\textbf{REPLICATED (qualitative, headline exercised).} The paper's central
phenomenological claim --- that the QAOA reachability deficit
\(f(p,\alpha,n) = \min_{\vec\gamma,\vec\beta}\langle\psi_p|H_{\rm SAT}|\psi_p\rangle - E_{\min}\)
grows monotonically with clause density \(\alpha=m/n\) and shrinks
(but does not vanish) with circuit depth \(p\) --- was independently
regenerated end-to-end on freshly sampled random 3-SAT instances at
\(n=6\), across \(\alpha\in\{0.5,1,2,3,5,7,10\}\) and
\(p\in\{1,2,4\}\), with a strictly monotone trend in every row and
column of the sweep and zero significant non-monotonicities.

\section{What was independently regenerated (headline-exercised rule)}
\begin{itemize}
  \item \textbf{Full QAOA statevector engine} written from scratch in NumPy
        (\(\sim\)150 lines, no Qiskit/PennyLane), verified by three unit
        tests (clause truth table, plus-state expectation, mixer unitarity).
  \item \textbf{Instance ensemble} regenerated: 15 random 3-SAT instances
        per \(\alpha\) (paper: 100), fresh seed \texttt{20260703}, canonical
        random-3-SAT ensemble (3 distinct vars per clause, each literal
        negated with prob.~\(\tfrac12\)).
  \item \textbf{Optimization} independently: COBYLA, 4 random restarts per
        (instance,\(p\)), \texttt{maxiter=300}, \texttt{rhobeg=0.3}.
  \item \textbf{Reachability-deficit curves} recomputed for each
        \((p,\alpha)\) cell and plotted vs.~\(\alpha\) (Fig.~1 analog).
\end{itemize}
This satisfies the headline-exercised standard: we did not merely quote
the paper's Figure 1 numbers, we \emph{regenerated} deficit curves from
fresh instances with our own optimizer and compared shape/ordering.

\section{Results summary}
\begin{center}\small
\begin{tabular}{cccccc}
\toprule
\(\alpha\) & \(m\) & \(p{=}1\)~\(f_{\rm mean}\pm\)SEM & \(p{=}2\) & \(p{=}4\) & \(\langle\min H_{\rm SAT}\rangle\)\\
\midrule
0.5 & 3  & 0.069\(\pm\)0.023 & 0.013\(\pm\)0.003 & 0.005\(\pm\)0.001 & 0.00\\
1.0 & 6  & 0.147\(\pm\)0.009 & 0.072\(\pm\)0.013 & 0.046\(\pm\)0.019 & 0.00\\
2.0 & 12 & 0.608\(\pm\)0.052 & 0.376\(\pm\)0.034 & 0.245\(\pm\)0.037 & 0.00\\
3.0 & 18 & 1.076\(\pm\)0.076 & 0.822\(\pm\)0.067 & 0.619\(\pm\)0.077 & 0.00\\
5.0 & 30 & 1.811\(\pm\)0.160 & 1.590\(\pm\)0.201 & 1.340\(\pm\)0.169 & 0.47\\
7.0 & 42 & 2.475\(\pm\)0.208 & 1.851\(\pm\)0.221 & 1.853\(\pm\)0.204 & 1.13\\
10.0& 60 & 2.783\(\pm\)0.216 & 2.561\(\pm\)0.208 & 2.140\(\pm\)0.206 & 2.73\\
\bottomrule
\end{tabular}
\end{center}
Monotonicity across \(\alpha\): 6/6 up-steps at every \(p\), zero
reversals. Monotonicity across \(p\) at each \(\alpha\): 7/7.
Deficit ratio \(f(\alpha{=}10)/f(\alpha{=}0.5)\) is 40.2 (p=1),
200.9 (p=2), 413.5 (p=4).

\section{Honest critique / limits of this replication}

\subsection{Density transition was reproduced only \emph{qualitatively}}
The paper's Fig.~1 shows a sharp knee near \(\alpha\approx 1\); our
data show the same knee (mean \(f\) at \(\alpha=0.5\) is \(\sim 10^{-2}\)
and jumps to order-1 by \(\alpha=3\)). However, we did \emph{not}
attempt a quantitative fit to a critical-density functional form
(e.g.\ finite-size scaling in \(\alpha_c(n)\)) --- the paper itself
frames the transition qualitatively at this \(n\), so this is
consistent with the paper's presentation but leaves the sharpness
of the transition uncalibrated.

\subsection{Circuit depth is 4$\times$ smaller than paper's ($p\le 4$ vs $p\in\{15,25,35\}$)}
This is the single largest methodological gap. Consequences:
\begin{itemize}
  \item Our absolute \(f\) values are systematically larger than the
        paper's (e.g.\ \(f(\alpha{=}10,p{=}4)\approx 2.14\) vs paper's
        \(\sim 1.5\) at \(p{=}15\)). Since \(f\) is monotone
        non-increasing in \(p\), our numbers are an upper bound on the
        paper's --- they cannot disprove the trend but they cannot
        confirm the exact numeric depth of the paper's curves either.
  \item Claim C3 (critical depth \(p^*\) for \(\eta\ge 0.95\) grows with
        \(\alpha\)) is \textbf{not tested} --- it needs \(p\) up to
        \(\sim 40\), which would be \(\sim 10\times\) our compute
        budget for this wave.
\end{itemize}

\subsection{Small instance count inflates SEM}
15 instances per cell (paper: 100) inflates our SEM by
\(\sqrt{100/15}\approx 2.6\). Worst signal-to-noise (\(\alpha{=}1,
p{=}4\): \(f_{\rm mean}/\)SEM \(\approx 2.4\)) is above 2\(\sigma\)
detection but not comfortable; the large-\(\alpha\) signal is
\(\gtrsim 10\sigma\).

\subsection{Small variable count $n=6$; no finite-size trend}
Paper uses \(n\in\{6,8,10,12\}\) to argue the deficit persists as \(n\)
grows. We ran \emph{only} \(n=6\). This means our replication does not
address the finite-size question at all --- we confirm the phenomenon
exists at the smallest \(n\) the paper reports and take on faith the
finite-size trend the paper claims.

\subsection{Secondary claims not tested}
\begin{itemize}
  \item \textbf{C2} (projector-mixer variant, paper's Fig.~2) --- not run.
  \item \textbf{C3} (\(p^*\) vs \(\alpha\), Fig.~3) --- not run (depth budget).
  \item \textbf{C4} (variational Grover \(p^*\sim\sqrt{N}\), Fig.~4) --- not run.
\end{itemize}
Only C1 (the paper's headline) is exercised.

\subsection{Comparison against alternative QAOA analyses not attempted}
We did not cross-check whether the deficits we observe are (i) an
optimization-landscape artifact (barren plateaus, concentration of
measure --- McClean et al.\ 2018, Wang et al.\ 2021), (ii) an expressibility
artifact (the standard X-mixer ansatz manifold really does not
contain the ground state at low \(p\), Larocca et al.\ 2022), or
(iii) genuinely both. The paper argues (ii); we did not test whether
our COBYLA/4-restart optimizer might have gotten stuck in local
minima and inflated \(f\) beyond the true reachability minimum. A
tighter test would repeat the sweep with L-BFGS-B, more restarts, or
BOBYQA and check that \(f\) is stable.

\subsection{No comparison against modern QAOA variants}
Since 2019 the literature has produced ADAPT-QAOA, warm-start QAOA,
multi-angle QAOA, and RQAOA. None of these were tested against the
same instance ensemble --- doing so would answer whether the
reachability deficit is a property of the \emph{ansatz family} or of
the underlying problem class.

\section{Reproducibility}
All code, seeds, per-instance JSON, CSV summary, figure, and full
515\,s sweep log are preserved under
\texttt{code/}, \texttt{data/}, \texttt{figures/}, \texttt{logs/},
and mirrored in \texttt{report/evidence/}. Reproduce with:
\begin{verbatim}
python3 -m venv .venv && source .venv/bin/activate
pip install numpy scipy matplotlib
python code/smoke.py                 # 3 unit tests (~1s)
python code/qaoa_reachability.py     # main sweep (~8-9 min CPU)
python code/plot_results.py          # figure + CSV
\end{verbatim}

\section{Bottom line}
The paper's headline claim survives independent replication:
QAOA \emph{does} exhibit an \(\alpha\)-dependent reachability deficit
that shrinks with \(p\) but persists at shallow depth. What our
replication does \emph{not} do is pin down the quantitative
critical-density behavior, test the deeper-\(p\) regime the paper
actually studies, or distinguish optimizer failure from ansatz
inexpressibility. The open-questions section enumerates concrete
next probes.

\end{document}
